\chapter{Introduction méthodologique}

Cette partie présente trois expériences pédagogiques significatives qui illustrent différentes facettes de mes compétences professionnelles développées dans le cadre de l'ingénierie de formation. Ces expériences ont été sélectionnées pour leur représentativité par rapport aux compétences du référentiel RNCP38152 du Master MEEF PIF \parencite{rncp38152}.

Chaque expérience est analysée selon une structure rigoureuse en quatre axes, conformément aux recommandations méthodologiques de l'INSPÉ de Normandie Caen / Unicaen \parencite{vademecum_vae_inspe}:

\begin{enumerate}[leftmargin=*]
    \item \textbf{Axe 1 - Contexte et intentions pédagogiques} : présentation du public, des objectifs d'apprentissage et des choix didactiques
    \item \textbf{Axe 2 - Déroulement et rôle spécifique} : description détaillée des actions menées et de ma posture professionnelle
    \item \textbf{Axe 3 - Résultats et évaluation} : analyse des productions des apprenants et de l'efficacité des dispositifs
    \item \textbf{Axe 4 - Analyse rétrospective et transférabilité} : réflexion critique sur les pratiques et correspondance avec le référentiel MEEF
\end{enumerate}

Cette approche structurée permet de démontrer non seulement la maîtrise technique et pédagogique, mais aussi la capacité de pratique réflexive, compétence centrale du métier d'ingénieur de formation et de formateur \parencite{perrenoud_pratique_reflexive,schon_reflective_practitioner}.

\section*{Cadre de référence et articulation des deux référentiels}

Pour faciliter la lecture par les membres du jury, les analyses d'expérience font référence conjointement :
\begin{itemize}[leftmargin=*]
    \item La \textbf{Maquette locale du Master MEEF PIF (Parcours FFMS)} de l'INSPÉ de Normandie Caen / Unicaen, structurée en \textbf{6 Unités d'Enseignement (UE1 à UE6)} et Éléments Constitutifs (EC).
    \item Le \textbf{Référentiel National de Certification RNCP38152} du Master MEEF, structuré en \textbf{9 Blocs de Compétences (BC01 à BC09)}.
\end{itemize}

Afin d'éviter toute ambiguïté, chaque compétence analysée dans les chapitres 9, 10 et 11 fait l'objet d'un \textbf{double étiquetage systématique} : \texttt{[UE/EC Unicaen | RNCP BC]}.

\begin{table}[H]
\centering
\caption{Table de correspondance unifiée entre la Maquette FFMS Unicaen et la certification RNCP38152}
\label{tab-correspondance-referentiels}
\footnotesize
\begin{tabular}{|p{7.5cm}|p{7.5cm}|}
\hline
\textbf{Maquette Unicaen MEEF PIF (Parcours FFMS)} & \textbf{Certification Nationale RNCP38152 (9 Blocs)} \\
\hline
\textbf{UE1 (EC 111-212)} : Cadres d'exercice et principes éthiques & \textbf{BC01} (Cadre éthique) \& \textbf{BC02} (Postures professionnelles) \\
\hline
\textbf{UE2 (EC 121-225)} : Conception, mise en œuvre et évaluation de projets & \textbf{BC03} (Ingénierie pédagogique) \& \textbf{BC04} (Conduite de projet) \\
\hline
\textbf{UE3 (EC 131-233)} : Les contextes et les publics du métier & \textbf{BC07} (Inclusion \& EBEP) \& \textbf{BC09} (Communauté éducative) \\
\hline
\textbf{UE4 (EC 141-242)} : Formation à et par la recherche & \textbf{BC08} (Recherche, veille et innovation) \\
\hline
\textbf{UE5 (EC 151-254)} : Analyse de l'activité professionnelle & \textbf{BC06} (Praticien réflexif et analyse de pratiques) \\
\hline
\textbf{UE6 (EC 161-263)} : Communication et numérique éducatif & \textbf{BC05} (Didactique \& Numérique) \& \textbf{BC09} (Animation du collectif) \\
\hline
\end{tabular}
\end{table}

\section*{Continuum d'ingénierie de formation des trois expériences}

Les trois expériences s'articulent selon une progression d'ingénierie pédagogique et de formation continue :
\begin{enumerate}[leftmargin=*]
    \item \textbf{Chapitre 9 (Expérience 1 - Atelier IA/\LaTeX{} Forgéales)} : L'ingénierie de la formation d'adultes et de pairs formateurs (dispositif hybride synchrone/asynchrone, posture de formateur et régulation réflexive).
    \item \textbf{Chapitre 10 (Expérience 2 - PCI Jeu Vidéo Interdisciplinaire)} : L'ingénierie de dispositifs complexes de pédagogie par projet (scénarisation, découplage des artefacts didactiques, travail collaboratif).
    \item \textbf{Chapitre 11 (Expérience 3 - Notation à l'Enveloppe)} : L'ingénierie de l'évaluation régulatrice et métacognitive (réponse directe aux défis d'évaluation du travail de groupe identifiés dans le Chapitre 10).
\end{enumerate}
